\section{Recordatorio de ingeniería: entre más estados, más se necesitan límites explícitos}
La capacidad de ejecución asíncrona del Agentic Loop depende, en esencia, de una máquina de estados. Si la definición de los estados no es clara o las condiciones de transición son ambiguas, el sistema cae fácilmente en bloqueos, ejecuciones repetidas o una recuperación incorrecta de recursos a lo largo de las llamadas a herramientas en múltiples rondas.

\begin{warningbox}{Riesgos comunes en la gestión de estados}
\begin{itemize}
  \item El sample tarda demasiado en llegar a Done, lo que provoca fugas en la cola
  \item Después de que la herramienta responde, no se regresa correctamente a Generating
  \item Hay demasiados parámetros de configuración, pero la lógica de transición no tiene un punto de entrada unificado
\end{itemize}
\end{warningbox}

\subsection{Resumen de la sección}
Lo que más teme un sistema asíncrono es que «los estados no estén claros». Definir con precisión los límites de los estados y las condiciones de transición es más importante que perseguir sin criterio una mayor concurrencia.

\newpage
